\documentclass[book.tex]{subfiles}

\begin{document}

\chapter{Deep Networks Tend to Gaussian Processes}
\label{chap:ntk}
\noindent\rule{11cm}{0.4pt} \\
\textbf{Prerequisites:}  \autoref{chap:ml_basics}, \autoref{chap:grad_descent}  \\
\textbf{Difficulty Level:} *** \\
\noindent\rule{11cm}{0.4pt}
\vspace{5mm}

%\cite{lee2017deep}

Consider a fully connected neural network with $L$ hidden layers. Let the $\ell$-th layer have $N_\ell$ hidden units, bias $b^\ell$ and nonlinearity $\phi$. $x \in \mathbb{R}^d$ is the input and $z^{\ell}$ is the output of the $\ell$-th layer. One of the most important mathematical theorems in deep learning proves that as $N_\ell \to \infty$, the deep neural network converges to a Gaussian process. 

\section{A Single Hidden Layer Network is a Gaussian Process}
The $i$-th component output $z_i$ of a single hidden layer network is given by 
\begin{align}
a_j &= \sigma \left ( b^0_j + \sum_{k=1}^d W^0_{jk} x_k \right ) \\
z_i &= b^1_i + \sum_{j=1}^{N} W^1_{ij} a_j 
\end{align}

Assume that weight parameters $W^0_{ij}, W^1_{ij}$ and bias parameters $b^0_j, b^1_j$ are given by independent, identically distributed random draws from a Gaussian distribution. We can treat inputs $x_k$ as constants. Then note that the sum term
\begin{align}
b^0_j + \sum_{k=1}^d W^0_{jk} x_k
\end{align}
itself follows a Gaussian distribution, and $a_j$ is drawn from the nonlinear transformation $\sigma$ of this Gaussian. Then it follows that $a_j$ and $a_{j'}$ are independent and identically distributed when $j \neq j'$ since the summation terms draw from independent distributions (note that $x_k$ are constant terms so it does not matter that they are shared.) 

The output term $z_i$ is then the sum of i.i.d terms. As $N \to \infty$, it follows from the central limit theorem that $z_i$ will converge to a Gaussian distribution (we can ignore the bias term as $N \to \infty$).

It follows then that 
\begin{align}
z^1_i \sim \mathcal{GP}(\mu^1, K^1)
\end{align} 
is given by a Gaussian process since any collection of these terms for different $i$ forms a multivariate Gaussian distribution. Assuming that random parameters have mean $0$, then, $\mu^1 = 0$. The covariance is given by

\begin{align}
K(x, x') &= \mathbb{E}[z_i(x)z_i(x')] \\
\end{align}

\section{Deep Neural Networks are Gaussian Processes}

An argument by induction can prove that as we add additional layers, a deeper neural network still converges to a Gaussian process as the width goes to infinity.
\end{document}